\chapter{The Designation Mode}
\label{ch:designation}

\begin{versebox}
One teaching, many designations ---\\
the Blessed One frames it fresh each time.\\
Same truth, a hundred angles:\\
every framing fits the same design.
\end{versebox}

\noindent Imagine you are a radiologist. A patient comes in with a complaint, and you have at your disposal not one but twenty imaging modalities --- X-ray, MRI, CT scan, ultrasound, PET scan, and a dozen more. Each one shows the same body. But each reveals different features. The X-ray shows bone; the MRI shows soft tissue; the PET scan shows metabolic activity. Same patient, same complaint --- but each image gives you a different angle on the diagnosis.\footnote{Pali: \textit{``Tattha katamo paññattihāro? `Ekaṃ bhagavā dhammaṃ paññattīhi vividhāhi desetī'\,''ti.} The mnemonic fragment is ``The Blessed One teaches one teaching through various designations'' (\textit{ekaṃ bhagavā dhammaṃ paññattīhi vividhāhi deseti}). The word \textit{paññatti} means ``designation, description, concept, making known.'' The Designation Mode teaches that the Buddha used many types of \textit{paññatti} --- origin-designation, comprehension-designation, abandonment-designation, development-designation, realization-designation, and at least fifteen more --- to frame the same teaching from different diagnostic angles.}

The Designation Mode says: the Buddha did the same thing with his teaching. He had one insight --- the Four Noble Truths --- and he designated it, framed it, labeled it in dozens of different ways, depending on what the listener needed to hear.

The simplest designation is the setting-down: ``This is suffering.'' When the Buddha said those words, he was laying down the First Noble Truth through the five aggregates, the six elements, the eighteen elements, the twelve sense-bases, the ten faculties. Same truth. Different frameworks. Different imaging modalities for the same diagnosis.\footnote{Pali: \textit{``Yā pakatikathāya desanā. Ayaṃ nikkhepapaññatti. Kā ca pakatikathāya desanā, cattāri saccāni. Yathā bhagavā āha `idaṃ dukkha'\,''nti ayaṃ paññatti pañcannaṃ khandhānaṃ channaṃ dhātūnaṃ aṭṭhārasannaṃ dhātūnaṃ dvādasannaṃ āyatanānaṃ dasannaṃ indriyānaṃ nikkhepapaññatti.''} The setting-down designation (\textit{nikkhepapaññatti}) is the most basic type: a direct placement of a truth into a categorical framework. ``This is suffering'' is the setting-down of the First Noble Truth through five aggregates, six elements, eighteen elements, twelve sense-bases, and ten faculties. One sentence, five frameworks.}

But this is only the beginning. The Netti identifies at least twenty types of designation --- twenty different ways to label what a passage is doing. And once you learn to see them, every sutta becomes a prism.

\section*{The Four Nutrients}

Here is how it works in practice. Consider the Buddha's teaching on the four nutrients:\footnote{Pali: \textit{``Kabaḷīkāre ce, bhikkhave, āhāre atthi rāgo atthi nandī atthi taṇhā, patiṭṭhitaṃ tattha viññāṇaṃ virūḷhaṃ. Yattha patiṭṭhitaṃ viññāṇaṃ virūḷhaṃ, atthi tattha nāmarūpassa avakkanti. Yattha atthi nāmarūpassa avakkanti, atthi tattha saṅkhārānaṃ vuddhi. Yattha atthi saṅkhārānaṃ vuddhi, atthi tattha āyatiṃ punabbhavābhinibbatti.''} From the Āhāra Sutta (SN 12.64). The four nutrients (\textit{āhāra}) are: (1) physical food (\textit{kabaḷīkāra}), (2) contact (\textit{phassa}), (3) mental intention (\textit{manosañcetanā}), (4) consciousness (\textit{viññāṇa}). The Netti quotes the full passage for physical food and consciousness, abbreviating contact and intention with ``\ldots and so on'' (\textit{pe}).}

\begin{quote}
If there is lust, monks, if there is delight, if there is craving for physical food --- consciousness becomes established there and grows. Where consciousness becomes established and grows, name-and-form descends. Where name-and-form descends, formations grow. Where formations grow, future rebirth is produced. Where future rebirth is produced, future birth, aging, and death arise. Where future birth, aging, and death arise --- that, monks, I say is accompanied by sorrow, by anguish, by despair.
\end{quote}

Read that passage once, and you see a chain of consequences: craving leads to consciousness, consciousness leads to becoming, becoming leads to birth, birth leads to suffering. A domino chain.

But the Designation Mode reads it twice. The first reading labels the passage as an \textit{origin-designation} --- it designates where suffering and its origin come from.\footnote{Pali: \textit{``Ayaṃ pabhavapaññatti dukkhassa ca samudayassa ca.''} ``This is the origin-designation (\textit{pabhavapaññatti}) of suffering and of the origin [of suffering].'' The passage, read forward (with craving present), reveals the mechanism by which suffering arises. This is the diagnostic angle.}

Now read the reverse:

\begin{quote}
If there is no lust, monks, if there is no delight, if there is no craving for physical food --- consciousness does not become established there, does not grow. Where consciousness does not become established and does not grow, name-and-form does not descend. Where name-and-form does not descend, formations do not grow \ldots future birth, aging, and death do not arise. That, monks, I say is without sorrow, without anguish, without despair.
\end{quote}

Same passage, reversed. But now the Designation Mode assigns not one label but \textit{four}. The reversed passage is simultaneously a comprehension-designation of suffering (you now understand what suffering really is), an abandonment-designation of the origin (you see what needs to be given up), a development-designation of the path (you see what needs to be cultivated), and a realization-designation of cessation (you see what can be attained).\footnote{Pali: \textit{``Ayaṃ pariññāpaññatti dukkhassa, pahānapaññatti samudayassa, bhāvanāpaññatti maggassa, sacchikiriyāpaññatti nirodhassa.''} Four designations in a single reversed passage: comprehension-designation (\textit{pariññāpaññatti}) of suffering, abandonment-designation (\textit{pahānapaññatti}) of the origin, development-designation (\textit{bhāvanāpaññatti}) of the path, realization-designation (\textit{sacchikiriyāpaññatti}) of cessation. The same nutriment teaching, read in reverse, reveals all Four Noble Truths simultaneously. This is what makes the Designation Mode so powerful: it shows you how the Buddha encoded multiple truths into a single passage.}

One teaching. Five designations. Read forward, it tells you how suffering arises. Read backward, it tells you --- in a single breath --- how suffering is understood, its origin abandoned, the path developed, and cessation realized.

\section*{Every Sense Is Impermanent}

The Netti demonstrates the same technique with another well-known sutta:

\begin{quote}
Develop concentration, monks. A concentrated monk, heedful, alert, and mindful, understands things as they really are. And what does he understand? That the eye is impermanent. That visible forms are impermanent. That eye-consciousness is impermanent. That eye-contact is impermanent. And whatever feeling arises dependent on eye-contact --- whether pleasant, painful, or neutral --- that too is impermanent.
\end{quote}

And so on through the ear, the nose, the tongue, the body, and the mind --- all six senses, all their objects, all their forms of consciousness, all their contacts, all their feelings. Every last one: impermanent.\footnote{Pali: \textit{``Samādhiṃ, bhikkhave, bhāvetha \ldots `cakkhu anicca'\,''nti yathābhūtaṃ pajānāti. `Rūpā aniccā'\,''ti yathābhūtaṃ pajānāti \ldots `mano anicco'\,''ti yathābhūtaṃ pajānāti \ldots tampi aniccanti yathābhūtaṃ pajānāti.''} From SN 35.99. The Netti abbreviates the middle four senses with \textit{pe} (``and so on''), quoting the eye and mind passages in full. The structure is: for each of the six senses, the concentrated monk understands five things as impermanent --- the sense-organ, its objects, its consciousness, its contact, and the feeling born from that contact. Six senses times five = thirty aspects of experience, all seen as impermanent.}

The Designation Mode labels this passage with four designations, in a specific order: first a development-designation (this is about developing the path --- concentration), then a comprehension-designation (this is about fully understanding suffering --- everything seen is impermanent), then an abandonment-designation (this abandons the origin --- the illusion of permanence), then a realization-designation (this realizes cessation --- the meditator ``sees things as they really are'').\footnote{Pali: \textit{``Ayaṃ bhāvanāpaññatti maggassa, pariññāpaññatti dukkhassa, pahānapaññatti samudayassa, sacchikiriyāpaññatti nirodhassa.''} Note the order differs from the nutriment passage: here the development-designation comes first, because the sutta opens with an instruction to \textit{develop} concentration. The Designation Mode is sensitive to the emphasis of each passage --- it does not always assign the same labels in the same order.}

Notice the order is different from the nutriment passage. That is the point. Different suttas enter the Four Noble Truths through different doors. The Designation Mode tells you which door you are using.

\section*{Scatter It All}

And then there is the Buddha at his most startling. He turns to the monk Rādha and says:

\begin{quote}
Scatter form, Rādha. Demolish it. Destroy it. Play with it. Practice for the destruction of craving through wisdom. When craving is destroyed, suffering is destroyed. When suffering is destroyed --- nirvana.
\end{quote}

And the same for feeling, perception, formations, and consciousness. Scatter them all.\footnote{Pali: \textit{``Rūpaṃ, rādha, vikiratha vidhamatha viddhaṃsetha vikīḷaniyaṃ karotha, paññāya taṇhakkhayāya paṭipajjatha. Taṇhakkhayā dukkhakkhayo, dukkhakkhayā nibbānaṃ. Vedanaṃ \ldots saññaṃ \ldots saṅkhāre \ldots viññāṇaṃ vikiratha vidhamatha viddhaṃsetha vikīḷaniyaṃ karotha, paññāya taṇhakkhayāya paṭipajjatha. Taṇhakkhayā dukkhakkhayo, dukkhakkhayā nibbānaṃ.''} From SN 23.2 (Satta Sutta). Four verbs in rapid succession: scatter (\textit{vikiratha}), demolish (\textit{vidhamatha}), destroy (\textit{viddhaṃsetha}), play with (\textit{vikīḷaniyaṃ karotha}). The intensity is remarkable --- the Buddha is telling Rādha to treat the five aggregates not as precious possessions but as things to be dismantled and examined, the way a child takes apart a toy. ``Play with it'' (\textit{vikīḷaniyaṃ karotha}) is particularly striking: don't fear the aggregates, don't cling to them --- play with them, see through them.}

This is language designed to shock. You are not supposed to preserve the aggregates, guard them, cherish them. You are supposed to pull them apart --- dismantle the very building blocks of experience until there is nothing left to cling to.

The Designation Mode assigns this passage a cascade of labels: it is a cessation-designation of nirvana, a disenchantment-designation of gratification, a comprehension-designation of suffering, an abandonment-designation of the origin, a development-designation of the path, and a realization-designation of cessation.\footnote{Pali: \textit{``Ayaṃ nirodhapaññatti nirodhassa, nibbidāpaññatti assādassa, pariññāpaññatti dukkhassa, pahānapaññatti samudayassa, bhāvanāpaññatti maggassa, sacchikiriyāpaññatti nirodhassa.''} Six designations for a single passage --- the most we have seen so far. The cessation-designation (\textit{nirodhapaññatti}) and disenchantment-designation (\textit{nibbidāpaññatti}) are new types, added to the four we have already met. The Netti's catalogue of designation-types keeps growing, because the more angles you examine a teaching from, the more you see.}

Six labels for one short passage. Each reveals a different aspect of what the Buddha is doing. He is not just telling Rādha to smash things. He is, simultaneously, showing him what cessation looks like, cultivating disenchantment, comprehending suffering, abandoning its origin, developing the path, and realizing the goal.

\section*{Three Turnings of the Wheel}

The most famous sutta in the entire Canon --- the very first sermon, the Dhammacakkappavattana Sutta --- gets the full treatment. The Buddha's announcement of the Four Noble Truths came in three waves, three turnings of the wheel. And each turning is a different type of designation.\footnote{Pali: \textit{``Idaṃ `dukkha'\,''nti me, bhikkhave, pubbe ananussutesu dhammesu cakkhuṃ udapādi, ñāṇaṃ udapādi, paññā udapādi, vijjā udapādi, āloko udapādi \ldots Taṃ kho panidaṃ dukkhaṃ pariññeyya\,''nti \ldots Taṃ kho panidaṃ dukkhaṃ pariññāta\,''nti \ldots''} The three turnings from SN 56.11. First turning: ``This is suffering'' (recognition). Second turning: ``This suffering should be fully understood'' (task). Third turning: ``This suffering has been fully understood'' (completion). The Netti maps each turning onto a different type of designation and a different type of wisdom.}

The first turning: ``This is suffering \ldots this is its origin \ldots this is its cessation \ldots this is the path.'' Straight declaration. The Designation Mode calls this a teaching-designation --- the Buddha is making the truths known for the first time, framing them as objects of hearing. This is the designation that plants wisdom-born-of-hearing in the listener, and it corresponds to the realization of the faculty called ``I shall know what is not yet known'' --- the first stirring of the spiritual life.\footnote{Pali: \textit{``Ayaṃ desanāpaññatti saccānaṃ, nikkhepapaññatti sutamayiyā paññāya, sacchikiriyāpaññatti anaññātaññassāmītindriyassa, pavattanāpaññatti dhammacakkassa.''} Four designations for the first turning alone: teaching-designation (\textit{desanāpaññatti}) of the truths, setting-down-designation of hearing-based wisdom (\textit{sutamayī paññā}), realization-designation of the faculty ``I shall know what is not yet known'' (\textit{anaññātaññassāmītindriya}), and turning-designation (\textit{pavattanāpaññatti}) of the Dhamma-wheel. This is the designation of the beginner: you hear the truth for the first time and something wakes up inside you.}

The second turning: ``This suffering should be fully understood \ldots this origin should be abandoned \ldots this cessation should be realized \ldots this path should be developed.'' Now the truths are framed not as facts to be heard but as tasks to be done. This is a development-designation --- the framing that plants wisdom-born-of-reflection, corresponding to the faculty called ``knowing.''\footnote{Pali: \textit{``Ayaṃ bhāvanāpaññatti maggassa, nikkhepapaññatti cintāmayiyā paññāya, sacchikiriyāpaññatti aññindriyassa.''} The second turning: development-designation of the path, setting-down of reflection-based wisdom (\textit{cintāmayī paññā}), realization of the faculty of ``knowing'' (\textit{aññindriya}). The shift from first to second turning is the shift from hearing to reflecting --- from ``I know what this is'' to ``I know what to do about it.''}

The third turning: ``This suffering has been fully understood \ldots this origin has been abandoned \ldots this cessation has been realized \ldots this path has been developed.'' Completed action. Done. This is the designation that plants wisdom-born-of-practice itself, corresponding to the faculty called ``one who has known'' --- the declaration of the arahant.\footnote{Pali: \textit{``Ayaṃ bhāvanāpaññatti maggassa, nikkhepapaññatti bhāvanāmayiyā paññāya, sacchikiriyāpaññatti aññātāvino indriyassa, pavattanāpaññatti dhammacakkassa.''} The third turning: development-designation of the path, setting-down of practice-based wisdom (\textit{bhāvanāmayī paññā}), realization of the faculty of ``one who has known'' (\textit{aññātāvindriya}), and turning-designation of the Dhamma-wheel. The three turnings correspond to the three types of wisdom (hearing → reflection → practice) and the three spiritual faculties (\textit{anaññātaññassāmītindriya} → \textit{aññindriya} → \textit{aññātāvindriya}). This mapping is one of the Netti's most elegant achievements.}

Three turnings. Three types of wisdom: hearing, reflection, practice. Three spiritual faculties: the awakening, the knowing, and the known. The same Four Noble Truths, but each time framed at a deeper level --- from recognition to task to completion. That is the Designation Mode in action: not a new teaching each time, but the same truth, re-designated for a different stage of understanding.

\section*{The Sage Who Broke Through}

The Netti then turns to a verse about the sage --- the awakened one who has crossed over:

\begin{versebox}
Having known both the measurable and the immeasurable,\\
the sage abandoned the formation of becoming.\\
Inwardly delighted, concentrated,\\
he broke through the armor like one self-born.
\end{versebox}

Every word gets its own designation. ``The measurable'' is the conditioned world --- all compounded things that can be measured, counted, weighed. ``The immeasurable'' is nirvana --- the unconditioned, beyond all quantity. ``Having known both'' is a direct-knowledge-designation: the sage knows the entire range of reality, from the lowest to the highest.\footnote{Pali: \textit{``Tulamatulañca sambhavaṃ, bhavasaṅkhāramavassaji muni; / Ajjhattarato samāhito, abhindi kavacamivattasambhava\,''nti.} (Cf.\ DN II.169.) \textit{``Tula\,''nti saṅkhāradhātu. `Atula\,''nti nibbānadhātu, `tulamatulañca sambhava\,''nti abhiññāpaññatti sabbadhammānaṃ. Nikkhepapaññatti dhammapaṭisambhidāya.''} ``The measurable'' (\textit{tula}) = the conditioned element (\textit{saṅkhāradhātu}). ``The immeasurable'' (\textit{atula}) = the nirvana-element (\textit{nibbānadhātu}). ``Having known both'' = direct-knowledge-designation (\textit{abhiññāpaññatti}) of all things, and setting-down-designation of analytical knowledge of the teaching (\textit{dhammapaṭisambhidā}).}

``Abandoned the formation of becoming'' --- that is a relinquishment-designation of the origin and a comprehension-designation of suffering. He gave up the engine that keeps the cycle turning.\footnote{Pali: \textit{``Bhavasaṅkhāramavassaji munī\,''ti pariccāgapaññatti samudayassa. Pariññāpaññatti dukkhassa.''} Relinquishment-designation (\textit{pariccāgapaññatti}) = the sage actively gave up the formations that produce becoming. Comprehension-designation (\textit{pariññāpaññatti}) = by giving them up, he fully understood what suffering is.}

``Inwardly delighted, concentrated'' --- a development-designation of mindfulness directed at the body, and a stability-designation of one-pointedness of mind. The sage is not straining. He is at home inside himself.\footnote{Pali: \textit{``Ajjhattarato samāhito\,''ti bhāvanāpaññatti kāyagatāya satiyā. Ṭhitipaññatti cittekaggatāya.''} Development-designation of body-based mindfulness (\textit{kāyagatā sati}); stability-designation (\textit{ṭhitipaññatti}) of mental one-pointedness (\textit{cittekaggatā}). The phrase ``inwardly delighted'' (\textit{ajjhattarato}) is glossed as mindfulness of the body --- delight turned inward, away from external sense-objects.}

``Broke through the armor like one self-born'' --- a deep-disenchantment-designation of the mind, an acquisition-designation of omniscience, and a breaking-through-designation of the shell of ignorance. The armor is the hard crust of delusion that encases every unawakened mind. The sage cracked it open, the way a chick breaks through its shell.\footnote{Pali: \textit{``Abhindi kavacamivattasambhava\,''nti abhinibbidāpaññatti cittassa, upādānapaññatti sabbaññutāya, padālanāpaññatti avijjaṇḍakosānaṃ.''} Three designations for a single metaphor: deep-disenchantment-designation (\textit{abhinibbidāpaññatti}) of the mind, acquisition-designation (\textit{upādānapaññatti}) of omniscience (\textit{sabbaññutā}), and breaking-through-designation (\textit{padālanāpaññatti}) of the eggshell of ignorance (\textit{avijjaṇḍakosa}). The image of breaking through an egg-shell of ignorance is vivid and ancient --- the sage is born twice, the second time by wisdom.}

Eight designations. From a four-line verse. Each one illuminates a different facet of what the sage accomplished --- and the Designation Mode keeps every facet visible at once.

\section*{Who Could Bend to Pleasure?}

Another verse. The Buddha asks a rhetorical question:

\begin{versebox}
Whoever has seen suffering and where it comes from ---\\
how could that person bend toward sense-pleasures?\\
Knowing that pleasures are the world's snare,\\
a mindful person trains for their removal.
\end{versebox}

The Netti takes this apart phrase by phrase, and the designations multiply.\footnote{Pali: \textit{``Yo dukkhamaddakkhi yatonidānaṃ, kāmesu so jantu kathaṃ nameyya; / Kāmā hi loke saṅgoti ñatvā, tesaṃ satīmā vinayāya sikkhe\,''ti.} Each phrase receives multiple designations: ``whoever saw suffering'' = synonym-designation + comprehension-designation of suffering; ``where it comes from'' = origin-designation + abandonment-designation of the origin; ``saw'' = synonym-designation of the knowledge-eye + penetration-designation; ``how could that person bend toward sense-pleasures'' = synonym-designation of sensual craving + commitment-designation; ``pleasures are the world's snare'' = seeing-from-the-enemy-angle designation (\textit{paccatthikato dassanapaññatti}) of pleasures; ``a mindful person'' = decumulation-designation for abandonment + setting-down of body-mindfulness + development-designation of the path; ``trains for their removal'' = penetration-designation of the removal of lust, hatred, and delusion.}

``Whoever has seen suffering'' --- two labels at once: a synonym-designation (``suffering'' is another word for the first truth) and a comprehension-designation (seeing it means understanding it). ``Where it comes from'' --- an origin-designation and an abandonment-designation. The word ``seen'' is itself a synonym-designation of the eye of knowledge and a penetration-designation.

``Pleasures are the world's snare'' --- here the Netti assigns a striking label: seeing-from-the-enemy-angle. Pleasures are not neutral. They are adversaries. And the traditional images are not gentle: pleasures are like a pit of hot coals, a slab of meat fought over by vultures, a blazing torch held into the wind, a pit with a venomous snake at the bottom.\footnote{Pali: \textit{``Kāmā hi aṅgārakāsūpamā maṃsapesūpamā pāvakakappā papātauragopamā ca.''} Four similes for sense-pleasures from the Potaliya Sutta (MN 54): (1) a pit of burning coals (\textit{aṅgārakāsu}), (2) a slab of meat (\textit{maṃsapesi}) --- over which dogs and vultures fight, (3) a blazing grass-torch (\textit{pāvaka}) held by someone walking into the wind, (4) a precipice with a venomous snake (\textit{papāta-uraga}). These are the ``enemy-angle'' designations: the teaching reframes pleasures from ``things I want'' to ``things that destroy me.''}

``A mindful person'' gets three designations: a decumulation-designation (mindfulness reduces the heap of defilements), a setting-down-designation of body-directed mindfulness, and a development-designation of the path.

And ``trains for their removal'' --- a penetration-designation of the removal of lust, hatred, and delusion. Not just any removal. Penetrative removal. The kind that goes all the way down.\footnote{Pali: \textit{``Vinayāya sikkhe\,''ti paṭivedhapaññatti rāgavinayassa dosavinayassa mohavinayassa.''} Penetration-designation (\textit{paṭivedhapaññatti}) of the removal (\textit{vinaya}) of the three root defilements: lust (\textit{rāga}), hatred (\textit{dosa}), delusion (\textit{moha}).}

And then the Netti goes further. The word ``person'' in the verse is itself a synonym-designation for the practitioner. And when that practitioner understands that pleasures are a snare, what does he do? He generates wholesome states that haven't yet arisen --- and that is an effort-designation of attaining what has not yet been attained. He maintains wholesome states that have already arisen --- and that is a heedfulness-designation of cultivation, a guarding-designation of wholesome qualities, and a stability-designation of the higher training in the mind.\footnote{Pali: \textit{``Jantū\,''ti vevacanapaññatti yogissa \ldots so kāmānaṃ anuppādāya kusale dhamme uppādayati \ldots Ayaṃ vāyāmapaññatti appattassa pattiyā. Nikkhepapaññatti oramattikāya asantuṭṭhiyā. Tattha so uppannānaṃ kusalānaṃ dhammānaṃ ṭhitiyā vāyamatī\,''ti ayaṃ appamādapaññatti bhāvanāya, nikkhepapaññatti vīriyindriyassa, ārakkhapaññatti kusalānaṃ dhammānaṃ, ṭhitipaññatti adhicittasikkhāya.''} A single word (``person,'' \textit{jantu}) triggers an entire chain of designations about the practitioner's activity: effort-designation (\textit{vāyāmapaññatti}), setting-down-designation of discontent with mere partial attainment (\textit{oramattikā asantuṭṭhi}), heedfulness-designation (\textit{appamādapaññatti}), setting-down of the energy-faculty (\textit{vīriyindriya}), guarding-designation (\textit{ārakkhapaññatti}) of wholesome qualities, and stability-designation (\textit{ṭhitipaññatti}) of the higher mind-training (\textit{adhicittasikkhā}).}

One verse. More than a dozen designations. Each one a window onto the same landscape.

\section*{Bound by Delusion}

One more verse, dark and striking:

\begin{versebox}
The world, bound by delusion, appears to be substantial ---\\
the fool, fettered to acquisitions, surrounded by darkness.\\
It seems as if it will last forever.\\
For one who sees, there is nothing there.
\end{versebox}

``Bound by delusion'' --- a teaching-designation of the distortions. The world is not actually substantial; it only looks that way because delusion has bound the viewer.\footnote{Pali: \textit{``Mohasambandhano loko, bhabbarūpova dissati; upadhibandhano bālo, tamasā parivārito; assirī viya khāyati, passato natthi kiñcana\,''nti.} (Cf.\ Ud 70.) \textit{``Mohasambandhano loko\,''ti desanāpaññatti vipallāsānaṃ.''} Teaching-designation (\textit{desanāpaññatti}) of the distortions (\textit{vipallāsa}). The entire world, as experienced by the deluded, is seen through the lens of the four distortions (permanence, pleasure, self, beauty). What appears ``substantial'' (\textit{bhabbarūpa}) is a mirage.}

``Appears to be substantial'' --- an inversion-designation. The world is designated as appearing in the opposite of its true nature.\footnote{Pali: \textit{``Bhabbarūpova dissatī\,''ti viparītapaññatti lokassa.''} Inversion-designation (\textit{viparītapaññatti}): the world appears inverted, substantial where it is insubstantial, permanent where it is fleeting.}

``The fool, fettered to acquisitions'' --- the Netti piles on labels: an origin-designation of evil desires, a function-designation of obsessions, a power-designation of defilements, a growth-designation of formations. The fool is trapped not by one chain but by many, each reinforcing the others.\footnote{Pali: \textit{``Upadhibandhano bālo\,''ti pabhavapaññatti pāpakānaṃ icchāvacarānaṃ, kiccapaññatti pariyuṭṭhānānaṃ. Balavapaññatti kilesānaṃ. Virūhanāpaññatti saṅkhārānaṃ.''} Four designations for the fool's bondage: origin-designation (\textit{pabhavapaññatti}) of evil desires (\textit{pāpaka icchāvacara}), function-designation (\textit{kiccapaññatti}) of obsessions (\textit{pariyuṭṭhāna}), power-designation (\textit{balavapaññatti}) of defilements (\textit{kilesa}), growth-designation (\textit{virūhanāpaññatti}) of formations (\textit{saṅkhāra}). The Designation Mode here reveals how a single phrase contains multiple diagnostic layers --- the fool's bondage is simultaneously analyzed from the angle of desire, of obsession, of defilement-power, and of karmic accumulation.}

``Surrounded by darkness'' --- a teaching-designation and a synonym-designation, because ``darkness'' is simply another name for ignorance.

``Seems as if it will last forever'' --- a seeing-designation through the divine eye, a setting-down-designation of the wisdom-eye. Two kinds of seeing, two kinds of designation.

And then the devastating last line: ``For one who sees, there is nothing there.'' The ``nothing'' is not nihilism. It is the penetration-designation of beings: when you see clearly, you see that there is no ``thing'' --- no lust, no hatred, no delusion --- remaining. The three defilements are ``something.'' Their absence is ``nothing.'' And that nothing is freedom.\footnote{Pali: \textit{``Passato natthi kiñcana\,''nti paṭivedhapaññatti sattānaṃ, rāgo kiñcanaṃ doso kiñcanaṃ moho kiñcanaṃ.''} Penetration-designation (\textit{paṭivedhapaññatti}) of beings. ``Something'' (\textit{kiñcana}) = lust, hatred, delusion. ``Nothing'' (\textit{natthi kiñcanaṃ}) = their absence. For the one who truly sees, the ``somethings'' have been removed, and what remains is ``nothing'' --- which is to say, everything that matters. A characteristically Buddhist inversion: emptiness is fullness, nothing is everything, the void is liberation.}

\section*{The Unborn}

The Designation Mode closes with one of the most famous passages in the entire Canon --- the Buddha's declaration of the unconditioned:\footnote{Pali: \textit{``Atthi, bhikkhave, ajātaṃ abhūtaṃ akataṃ asaṅkhataṃ, no cetaṃ, bhikkhave, abhavissa ajātaṃ abhūtaṃ akataṃ asaṅkhataṃ. Nayidha jātassa bhūtassa katassa saṅkhatassa nissaraṇaṃ paññāyetha. Yasmā ca kho, bhikkhave, atthi ajātaṃ abhūtaṃ akataṃ asaṅkhataṃ, tasmā jātassa bhūtassa katassa saṅkhatassa nissaraṇaṃ paññāyatī\,''ti.} From Ud 73 (Pāṭaligāmiya Vagga). This is one of the most philosophically significant passages in the entire Pali Canon. The argument is logical: \textit{if} there were no unborn, unbecome, uncreated, unconditioned, \textit{then} no escape from the born, become, created, conditioned could be discerned. \textit{But} since there \textit{is} an unborn, unbecome, uncreated, unconditioned, \textit{therefore} an escape from the born, become, created, conditioned \textit{is} discerned.}

\begin{quote}
There is, monks, an unborn, unbecome, uncreated, unconditioned. If there were no unborn, unbecome, uncreated, unconditioned, then no escape from the born, the become, the created, the conditioned could ever be discerned. But since there \textit{is} an unborn, unbecome, uncreated, unconditioned --- therefore an escape from the born, the become, the created, the conditioned \textit{is} discerned.
\end{quote}

The logic is as clean as a mathematical proof. If X did not exist, then escape from not-X would be impossible. But escape from not-X \textit{is} possible. Therefore X exists.

The Designation Mode applies its final layer of analysis. ``There is an unborn'' --- a teaching-designation and a synonym-designation of nirvana. ``If there were no unborn \ldots no escape could be discerned'' --- a synonym-designation of the conditioned and an analogy-designation: the conditioned is being used to point, by contrast, to the unconditioned. ``Since there is an unborn'' --- a synonym-designation and an illumination-designation of nirvana. And the clincher: ``Therefore an escape is discerned'' --- a synonym-designation of nirvana, a leading-out-designation of the path, and an escape-designation from the cycle of existence.\footnote{Pali: \textit{``No cetaṃ, bhikkhave, abhavissa ajātaṃ abhūtaṃ akataṃ asaṅkhata\,''nti desanāpaññatti nibbānassa vevacanapaññatti ca. `Nayidha jātassa bhūtassa katassa saṅkhatassa nissaraṇaṃ paññāyethā\,''ti vevacanapaññatti saṅkhatassa upanayanapaññatti ca. `Yasmā ca kho, bhikkhave, atthi ajātaṃ abhūtaṃ akataṃ asaṅkhata\,''nti vevacanapaññatti nibbānassa jotanāpaññatti ca. `Tasmā jātassa bhūtassa katassa saṅkhatassa nissaraṇaṃ paññāyatī\,''ti ayaṃ vevacanapaññatti nibbānassa, niyyānikapaññatti maggassa, nissaraṇapaññatti saṃsārato.''} The final analysis: teaching-designation (\textit{desanāpaññatti}) + synonym-designation (\textit{vevacanapaññatti}) of nirvana; synonym-designation + analogy-designation (\textit{upanayanapaññatti}) of the conditioned; synonym-designation + illumination-designation (\textit{jotanāpaññatti}) of nirvana; synonym-designation + leading-out-designation (\textit{niyyānikapaññatti}) of the path + escape-designation (\textit{nissaraṇapaññatti}) from the round of existence (\textit{saṃsāra}). The Netti reads the passage as a multi-layered map: every clause simultaneously describes nirvana, the conditioned world, the path, and the escape.}

\ornament

More than twenty types of designation. Passages that carry one label, or four, or six, or eight. Verses where every word opens a new window. And through every window: the same view. The Four Noble Truths, designated and re-designated in an inexhaustible variety of framings.

The point is not to memorize the labels. The point is to develop the eyes of a radiologist. When you read a sutta, do not just read what it says. Ask: what type of designation is this? Is it teaching me about the origin of suffering? The path? The goal? Is it framing the truth as something to be heard, or reflected on, or practiced? Is it showing me the enemy from the front, or revealing the escape from behind?

Every sentence the Buddha spoke was loaded with more meaning than a single reading can extract. The Designation Mode is the set of lenses that lets you see it all.\footnote{Pali: \textit{``Tenāha āyasmā mahākaccāyano `ekaṃ bhagavā dhammaṃ, paññattīhi vividhāhi desetī\,''ti.''} The closing formula: ``Therefore the Venerable Mahākaccāyana said: `The Blessed One teaches one teaching through various designations.'\,'' \textbf{A note on our translation}: the original Pali in this section is highly repetitive in structure --- each passage is followed by a formulaic list of designation-types. We have woven this material into a flowing narrative, varying the presentation for each example. The designation-types and their assignments are preserved faithfully in the endnotes. The overall structure --- one teaching, many framings --- is the Netti's central claim, and we have kept it front and center throughout.}
